% ===== PRÉAMBULE CANONIQUE — à coller VERBATIM en tête de CHAQUE .tex =====
\documentclass[11pt,a4paper]{article}
\usepackage[utf8]{inputenc}\usepackage[T1]{fontenc}\usepackage{lmodern}
\usepackage[french]{babel}
\usepackage{amsmath,amssymb,mathtools}\usepackage{icomma}
\usepackage[version=4]{mhchem}          % \ce{...} : équations & formules
\usepackage{siunitx}\sisetup{locale=FR} % \qty{}{}, \si{}, \ang{}
\usepackage{chemfig}                    % \chemfig{...} : structures développées
\usepackage{xcolor}\usepackage{colortbl}
\usepackage{tikz}\usepackage{pgfplots}\pgfplotsset{compat=1.18}
\usepackage[most]{tcolorbox}\usepackage{enumitem}\usepackage{array,booktabs}
\usepackage[a4paper,margin=2.2cm]{geometry}
\usetikzlibrary{arrows.meta,calc,angles,quotes,positioning,patterns,backgrounds,shapes.geometric}
\definecolor{primary}{RGB}{222,49,99}\definecolor{primarydark}{RGB}{150,22,55}
\definecolor{defcol}{RGB}{0,114,178}\definecolor{methcol}{RGB}{52,92,148}
\definecolor{excol}{RGB}{0,135,90}\definecolor{attncol}{RGB}{200,40,55}
\tcbset{boxbase/.style={enhanced,breakable,boxrule=0.9pt,arc=2.5pt,
  left=3mm,right=3mm,top=2.3mm,bottom=2.3mm,fonttitle=\bfseries,coltitle=white}}
% --- boîtes TUTORIEL (noms EXACTS, ne pas renommer) ---
\newtcolorbox{cle}[1][]{boxbase,colback=defcol!6,colframe=defcol,
  title={\textsc{À retenir}\ifx\relax#1\relax\relax\else\textmd{ : \itshape #1}\fi}}
\newtcolorbox{methode}[1][]{boxbase,colback=methcol!7,colframe=methcol,sharp corners,
  title={\textsc{Méthode}\ifx\relax#1\relax\relax\else\textmd{ --- \itshape #1}\fi}}
\newtcolorbox{exemplebox}[1][]{boxbase,colback=excol!5,colframe=excol,
  title={\textsc{Exemple}\ifx\relax#1\relax\relax\else\textmd{ : \itshape #1}\fi}}
\newtcolorbox{piege}[1][]{boxbase,colback=attncol!6,colframe=attncol,
  title={\textsc{Piège}\ifx\relax#1\relax\relax\else\textmd{ : \itshape #1}\fi}}
\newtcolorbox{remarque}[1][]{boxbase,colback=black!4,colframe=black!45,
  title={\textsc{Remarque}\ifx\relax#1\relax\relax\else\textmd{ : \itshape #1}\fi}}
% --- boîtes EXERCICE / CORRIGÉ (noms EXACTS, auto-comptées) ---
\newtcolorbox[auto counter]{exo}[1][]{boxbase,colback=primary!4,colframe=primary,
  title={\textsc{Exercice}~\thetcbcounter\ifx\relax#1\relax\relax\else\textmd{ --- \itshape #1}\fi}}
\newtcolorbox[auto counter]{corrige}[1][]{boxbase,colback=excol!4,colframe=excol,
  title={\textsc{Corrigé}~\thetcbcounter\ifx\relax#1\relax\relax\else\textmd{ --- \itshape #1}\fi}}
\newcommand{\R}{\mathbb{R}}\newcommand{\N}{\mathbb{N}}\newcommand{\Z}{\mathbb{Z}}\newcommand{\ds}{\displaystyle}
% (un \usepackage{fancyhdr}+en-tête est possible ; il est ignoré côté web)
% =========================================================================
\begin{document}

\begin{exo}[masse volumique hors CNTP]
On utilise $\rho=\dfrac{pM}{RT}$ avec $R=\num{0,082}$ (unités \si{atm}, \si{\litre}, \si{\kelvin}).
\begin{enumerate}[label=\alph*)]
  \item \ce{O2} ($M=\qty{32}{\gram\per\mole}$) sous $p=\qty{2.0}{atm}$ à $T=\qty{300}{\kelvin}$ :
        calcule $\rho$.
  \item \ce{N2} ($M=\qty{28}{\gram\per\mole}$) sous $p=\qty{1.0}{atm}$ à $T=\qty{250}{\kelvin}$ :
        calcule $\rho$.
\end{enumerate}
\end{exo}

\begin{exo}[identifier un gaz]
Détermine la masse molaire, puis propose un gaz.
\begin{enumerate}[label=\alph*)]
  \item Aux CNTP, un gaz a une masse volumique $\rho=\qty{1.43}{\gram\per\litre}$.
  \item Un gaz a une densité $d=1{,}53$ par rapport à l'air.
\end{enumerate}
\end{exo}

\begin{exo}[influence de la température sur $\rho$]
Aux CNTP, un gaz a une masse volumique $\rho_1=\qty{1.50}{\gram\per\litre}$ (à $T_1=\qty{273}{\kelvin}$).
On le chauffe à $T_2=\qty{410}{\kelvin}$ en maintenant la pression constante.
\begin{enumerate}[label=\alph*)]
  \item À $p$ et $M$ fixés, comment $\rho$ dépend-elle de $T$ ? Écris la relation entre $\rho_1$,
        $\rho_2$, $T_1$ et $T_2$.
  \item Calcule la nouvelle masse volumique $\rho_2$.
\end{enumerate}
\end{exo}

\begin{exo}[masse molaire par pesée d'un gaz]
On pèse $m=\qty{1.80}{\gram}$ d'un gaz qui occupe $V=\qty{1.0}{\litre}$ sous $p=\qty{1.0}{atm}$ à
$T=\qty{300}{\kelvin}$ ($R=\num{0,082}$).
\begin{enumerate}[label=\alph*)]
  \item Écris la formule donnant $M$ en fonction de $m$, $R$, $T$, $p$, $V$.
  \item Calcule $M$ et propose un gaz.
\end{enumerate}
\end{exo}

\end{document}
